\section{Usuarios con discapacidad y funcionalidad asociada}
\label{sec:discapacidad}

Los \textit{stakeholders} \sk{4}--\sk{7} forman parte del usuario nuclear
del simulador. La estrategia es ofrecer \textbf{equivalencia funcional}:
toda capacidad del candidato general (\sk{1}) debe estar disponible para
ellos por una ruta accesible. Se trabaja sobre WCAG 2.2 nivel AA, con
verificación automática en CI mediante \texttt{axe-core} y
\texttt{Lighthouse}, y validación manual con representantes de cada
perfil. Esta sección integra la respuesta del proveedor Claude
(prompt 02), seleccionada por su rigor en la cita de criterios WCAG
concretos y por proponer métricas sustitutas que evitan penalizar
injustamente al usuario.

\subsection{Discapacidad visual (\sk{4})}

\subsubsection*{A. Necesidades específicas}
\begin{itemize}
  \item Acceso completo al contenido textual y gráfico de la sala
        virtual mediante lector de pantalla (NVDA, JAWS, VoiceOver),
        incluyendo el estado del aura del entrevistador y los
        indicadores de progreso de la sesión.
  \item Sustitución del canal visual del aura por un canal sonoro o
        textual equivalente, ya que los degradados luminosos son
        inaccesibles para usuarios ciegos y de bajo contraste
        discriminable para usuarios con baja visión.
  \item Control de tipografía, contraste y zoom del cliente sin que se
        rompan los flujos de entrevista, conservando la persistencia
        de configuración entre sesiones.
  \item Eliminación de la métrica de contacto visual del cómputo del
        puntaje global, sustituida por una métrica de orientación
        auditiva equivalente.
  \item Información clara y anticipada sobre el uso de la cámara:
        posibilidad de mantenerla apagada sin perder funcionalidades
        evaluables.
\end{itemize}

\subsubsection*{B. Funcionalidades verificables}
\begin{enumerate}
  \item Operación al 100\,\% por teclado en todos los flujos
        (login, configuración, entrevista, revisión de feedback, peer
        mock), conforme al criterio WCAG 2.2 2.1.1 \textit{Keyboard};
        verificable mediante recorrido E2E con Playwright sin uso de
        puntero.
  \item Todo elemento interactivo expone nombre, rol y estado
        accesibles vía ARIA, conforme al criterio WCAG 2.2 4.1.2
        \textit{Name, Role, Value}; verificable con \texttt{axe-core}
        en CI con cero violaciones de severidad
        \textit{serious}/\textit{critical}.
  \item El aura luminosa se traduce en un canal sonoro paralelo
        (\textit{earcons} no verbales con mapeo frecuencia-fluidez y
        panorámica-ritmo) que se activa con un \textit{toggle}
        persistido en el perfil; verificable con E2E que valida la
        emisión del \textit{earcon} ante eventos sintéticos.
  \item La métrica de contacto visual se reemplaza automáticamente por
        \emph{orientación de voz hacia el micrófono} cuando el usuario
        declara discapacidad visual o desactiva la cámara; verificable
        por telemetría del cambio de métrica activa y por test
        unitario del módulo de \textit{scoring}.
  \item El cliente cumple WCAG 2.2 1.4.3 \textit{Contrast (Minimum)}
        con ratio mínimo de 4.5:1 en texto y 3:1 en componentes, y
        WCAG 2.2 1.4.4 \textit{Resize Text} permitiendo zoom hasta
        200\,\% sin pérdida de contenido ni funcionalidad; verificable
        con Lighthouse Accessibility (puntaje mínimo 95) y pruebas
        visuales con zoom forzado.
  \item El feedback final se entrega también como reporte estructurado
        en HTML semántico y descargable en texto plano, navegable por
        encabezados y regiones; verificable con \texttt{axe-core} sobre
        el reporte y con E2E que confirma jerarquía h1--h2--h3 sin
        saltos.
\end{enumerate}

\subsection{Discapacidad auditiva (\sk{5})}

\subsubsection*{A. Necesidades específicas}
\begin{itemize}
  \item Equivalente textual en tiempo real para toda emisión de voz del
        entrevistador LLM, con latencia compatible con la cadencia
        conversacional de una entrevista.
  \item Posibilidad de responder en lengua de señas o por escrito sin
        que se degrade el análisis de competencias comunicativas.
  \item Sustitución de las métricas dependientes de prosodia (tono,
        pausas, muletillas) por métricas observables en el canal
        alternativo elegido por el usuario.
  \item Indicadores visuales redundantes para cualquier alerta o
        evento que en el flujo estándar se señaliza por audio.
  \item Compatibilidad con auriculares con bobina inductiva (T-coil) y
        control independiente del volumen del entrevistador respecto al
        volumen de notificaciones.
\end{itemize}

\subsubsection*{B. Funcionalidades verificables}
\begin{enumerate}
  \item Subtítulos sincronizados de las preguntas del entrevistador
        mediante TTS con \textit{captions} WebVTT, latencia mediana
        $\le 500\,\mathrm{ms}$ \texttt{(referencial)}, conforme a
        WCAG 2.2 1.2.4 \textit{Captions (Live)}; verificable por
        telemetría de latencia y E2E que valida la presencia y
        sincronía de pistas WebVTT.
  \item Modo \emph{respuesta escrita}: el turno se cierra mediante
        envío de texto en lugar de detección de fin de habla;
        verificable con E2E que completa una entrevista íntegra en
        este modo y telemetría que registra el modo activo por sesión.
  \item Cuando se selecciona discapacidad auditiva o el modo escrito,
        las métricas de fluidez verbal y muletillas se sustituyen por
        \emph{coherencia y densidad léxica} sobre la transcripción del
        usuario; verificable con test unitario del módulo de
        \textit{scoring}.
  \item Todos los eventos de audio (turno asignado, alerta de tiempo,
        cambio de fase, desconexión en peer mock) tienen contraparte
        visual persistente de al menos $3\,\mathrm{s}$
        \texttt{(referencial)}, conforme a WCAG 2.2 1.4.13
        \textit{Content on Hover or Focus}; verificable con E2E que
        dispara eventos sintéticos y valida la aparición en el DOM.
  \item El componente de subtítulos cumple WCAG 2.2 1.4.12
        \textit{Text Spacing} y permite ajuste de tamaño, color de
        fondo y opacidad sin recargar la sesión; verificable con
        \texttt{axe-core} y E2E de persistencia entre sesiones.
  \item La sala de peer mock incluye un canal de chat textual paralelo
        al canal de voz, con historial accesible durante toda la
        sesión; verificable con E2E multiusuario que confirma entrega
        bidireccional en menos de $1\,\mathrm{s}$ \texttt{(referencial)}.
\end{enumerate}

\subsection{Discapacidad motora (\sk{6})}

\subsubsection*{A. Necesidades específicas}
\begin{itemize}
  \item Operación del simulador sin requerir gestos finos de mouse,
        soportando teclado, conmutadores (\textit{switch access}) y
        comandos de voz.
  \item Tolerancia a tiempos de respuesta variables: ningún flujo
        crítico debe expirar por inactividad sin posibilidad de
        extensión.
  \item Eliminación o sustitución de las métricas de postura y
        gesticulación, que penalizarían a usuarios con movilidad
        reducida del tronco, brazos o cabeza.
  \item Encuadre flexible de la cámara: aceptación de planos no
        convencionales (lateral, sentado en silla de ruedas, decúbito)
        sin que la detección considere la imagen como inválida.
  \item Áreas de activación amplias y sin requerimiento de pulsaciones
        simultáneas o sostenidas.
\end{itemize}

\subsubsection*{B. Funcionalidades verificables}
\begin{enumerate}
  \item Toda acción ejecutable mediante un único punto de entrada
        (teclado, \textit{switch} o voz) sin combinaciones simultáneas
        obligatorias, conforme a WCAG 2.2 2.1.4
        \textit{Character Key Shortcuts} y 2.5.7
        \textit{Dragging Movements}; verificable con E2E
        \textit{keyboard-only} y otro \textit{voice-only}.
  \item Todos los objetivos táctiles tienen tamaño mínimo de
        $24\times 24$\,px CSS, conforme al criterio WCAG 2.2 2.5.8
        \textit{Target Size (Minimum)}; verificable con auditoría
        automatizada que mide el \textit{bounding box} de todos los
        elementos interactivos.
  \item Los temporizadores de respuesta del entrevistador son
        configurables (sin límite, estándar, $\times 2$, $\times 3$) y
        notifican antes de expirar con opción de prórroga, conforme a
        WCAG 2.2 2.2.1 \textit{Timing Adjustable}; verificable con E2E
        y telemetría del temporizador efectivo por sesión.
  \item Cuando el usuario declara discapacidad motora, las métricas de
        postura y gesticulación se reemplazan por
        \emph{estabilidad del encuadre facial} (presencia y centrado
        relativo del rostro) sin penalizar inclinaciones, asimetrías
        ni ausencia de movimiento corporal; verificable con test
        unitario del \textit{pipeline} MediaPipe y regresión sobre
        \textit{datasets} sintéticos de planos no convencionales.
  \item Capa de comandos de voz para todos los controles
        (\textit{``siguiente pregunta''}, \textit{``pausar''},
        \textit{``repetir''}), independiente del canal de respuesta
        sustantiva, conforme a WCAG 2.2 2.5.4 \textit{Motion
        Actuation}; verificable con E2E que dispara los comandos por
        API sintética y valida los efectos en el DOM.
  \item El sistema permite recalibrar el encuadre de la cámara
        declarando explícitamente la posición del usuario (de frente,
        lateral, sentado bajo, recostado) sin que la detección marque
        la imagen como no válida; verificable con test de integración
        que inyecta \textit{frames} de cada posición.
\end{enumerate}

\subsection{Discapacidad cognitiva y neurodivergencia (\sk{7})}

\subsubsection*{A. Necesidades específicas}
\begin{itemize}
  \item Reducción configurable de estímulos visuales y sonoros: el
        aura, las animaciones de fondo y los efectos de transición
        pueden saturar a usuarios con TEA o TDAH.
  \item Lenguaje claro en preguntas y feedback, con posibilidad de
        reformular en versión simplificada sin alterar el contenido
        evaluativo.
  \item Soporte tipográfico para usuarios con dislexia (fuentes
        específicas, espaciado aumentado, alineación a la izquierda
        sin justificado).
  \item Previsibilidad y control sobre la duración y estructura de la
        sesión, evitando cambios de fase sin aviso.
  \item Posibilidad de pausar la sesión sin penalización ni pérdida de
        progreso, en respuesta a sobrecarga sensorial o de atención.
\end{itemize}

\subsubsection*{B. Funcionalidades verificables}
\begin{enumerate}
  \item Modo de \emph{estimulación reducida} que desactiva el aura, las
        animaciones de fondo, los \textit{earcons} no esenciales y las
        transiciones, dejando únicamente indicadores estáticos de
        estado; verificable con E2E que confirma la ausencia de
        animaciones CSS activas (\texttt{prefers-reduced-motion}
        honrado), conforme a WCAG 2.2 2.3.3 \textit{Animation from
        Interactions}, y con auditoría Lighthouse.
  \item El usuario puede limitar el aura a un máximo de dos métricas
        de su elección entre las disponibles; verificable con E2E que
        selecciona la configuración y valida que el componente del
        aura solo renderiza los canales seleccionados.
  \item Todas las preguntas del entrevistador disponen de una versión
        en lenguaje claro (oraciones cortas, voz activa, una idea por
        oración) accesible mediante un control persistente; verificable
        con test unitario sobre el módulo de generación y métrica de
        legibilidad automatizada (índice INFLESZ o similar) con puntaje
        mínimo \texttt{(referencial)}.
  \item Control de tipografía con opciones validadas para dislexia
        (fuente con \textit{tracking} aumentado y espaciado interlínea
        $\ge 1.5$), conforme a WCAG 2.2 1.4.12 \textit{Text Spacing};
        verificable con E2E que aplica la configuración y valida
        \texttt{line-height}, \texttt{letter-spacing} y
        \texttt{word-spacing} en el DOM.
  \item Pausa y reanudación en cualquier punto sin pérdida del estado
        conversacional ni penalización en el puntaje, con estado
        conservado por al menos $24\,\mathrm{h}$ \texttt{(referencial)},
        conforme a WCAG 2.2 2.2.1 \textit{Timing Adjustable} y 2.2.6
        \textit{Timeouts}; verificable con E2E que pausa, cierra el
        navegador, reabre y valida la continuidad del contexto.
  \item Antes de iniciar cada fase, el simulador anuncia su duración
        estimada, su estructura y los criterios que se evaluarán, en
        formato textual persistente, conforme a la práctica de
        previsibilidad asociada a WCAG 2.2 3.2.5 \textit{Change on
        Request}; verificable con E2E que confirma la presencia del
        bloque de anuncio en cada transición de fase.
\end{enumerate}

\subsection{Verificación}

La validación combina cuatro herramientas y cuatro hitos por perfil:

\paragraph{Herramientas} \texttt{axe-core} integrado en CI con umbral
cero para violaciones \textit{serious} y \textit{critical};
\texttt{Lighthouse} con puntaje Accessibility mínimo de 95;
\texttt{Playwright} para E2E que cubren los flujos
\textit{keyboard-only}, modo escrito, comandos de voz y configuración
de tipografía; panel de telemetría propio que registra modo activo,
métricas sustitutas, latencia de subtítulos y temporizadores efectivos.

\paragraph{Hitos}
\begin{enumerate}
  \item Revisión heurística experta antes de la primera prueba con
        usuarios.
  \item Prueba moderada con un usuario representativo por perfil
        (\sk{4}--\sk{7}) sobre una entrevista completa.
  \item Peer mock con al menos un par mixto que incluya un usuario
        del perfil correspondiente.
  \item Auditoría final externa con conformidad WCAG 2.2 nivel AA
        documentada en informe firmado, previa a la entrega.
\end{enumerate}
